% =====================================================================
%  Determinación de la velocidad del sonido mediante un tubo resonante
%  Procedimiento completo y análisis físico del Asistente de Resonancia
%  Compilar con: pdflatex Procedimiento_Resonancia.tex  (dos veces)
% =====================================================================
\documentclass[11pt,a4paper]{article}

\usepackage[utf8]{inputenc}
\usepackage[T1]{fontenc}
\usepackage{lmodern}
\usepackage[spanish,es-noshorthands]{babel}
\decimalpoint
\usepackage{amsmath,amssymb}
\usepackage{siunitx}
\sisetup{output-decimal-marker={.},separate-uncertainty=true}
\usepackage{booktabs}
\usepackage[margin=2.4cm]{geometry}
\usepackage{xcolor}
\usepackage{enumitem}
\usepackage{hyperref}
\hypersetup{colorlinks=true,linkcolor=blue!60!black,urlcolor=blue!60!black}

\newcommand{\vsound}{v_{\text{aire}}}
\newcommand{\vexp}{v_{\text{exp}}}
\newcommand{\vteo}{v_{\text{teo}}}

\title{\textbf{Determinación de la velocidad del sonido con un tubo resonante}\\[2mm]
\large Procedimiento experimental y fundamento físico del \emph{Asistente de Resonancia}}
\author{Práctica de Laboratorio de Ondas}
\date{\today}

\begin{document}
\maketitle

\begin{abstract}
\noindent Se determina la velocidad de propagación del sonido en el aire a partir de las
resonancias de una columna de aire de longitud variable, excitada por un diapasón de
frecuencia conocida. El documento detalla el fundamento físico, el procedimiento paso a paso
asistido por el programa, los dos métodos de cálculo (punto único y diferencia de modos por
regresión lineal), la propagación de incertidumbres y un ejemplo numérico completo para un
diapasón de \SI{1024}{\hertz}. Se discuten además todas las fuentes de error y por qué el
procedimiento es metrológicamente robusto.
\end{abstract}

\section{Fundamento físico}

\subsection{Ondas estacionarias en un tubo cerrado--abierto}
El tubo resonante está \textbf{cerrado} en su extremo inferior por la superficie del agua y
\textbf{abierto} en la boca superior, donde se coloca el diapasón. Estas condiciones de
contorno imponen, para la onda de \emph{desplazamiento} de las partículas de aire:
\begin{itemize}[nosep]
  \item un \textbf{nodo} de desplazamiento en la superficie del agua (el aire no puede moverse contra el líquido);
  \item un \textbf{vientre} de desplazamiento en la boca abierta (máxima libertad de movimiento).
\end{itemize}
La distancia entre un nodo y el vientre contiguo es un cuarto de longitud de onda, $\lambda/4$.
Por ello sólo resuenan las longitudes de columna de aire que cumplen
\begin{equation}
  L_m \;=\; m\,\frac{\lambda}{4} \;-\; e, \qquad m = 1,\,3,\,5,\,7,\dots
  \label{eq:resonancia}
\end{equation}
donde $m$ recorre únicamente los \textbf{armónicos impares} (modos 1, 3, 5\dots) y $e$ es la
corrección de extremo (Sección~\ref{sec:correccion}). Dos resonancias consecutivas están
separadas exactamente por
\begin{equation}
  \Delta L = L_{m+2}-L_m = \frac{\lambda}{2}.
  \label{eq:separacion}
\end{equation}

\subsection{Corrección de extremo}\label{sec:correccion}
El vientre de presión no se forma exactamente en el plano de la boca, sino una pequeña
distancia por fuera. Para un tubo cilíndrico con un solo extremo abierto, esa corrección vale
\begin{equation}
  e \;\approx\; 0.6\,r \;=\; 0.3\,D,
  \label{eq:endcorr}
\end{equation}
con $r$ el radio y $D$ el diámetro interno del tubo. La ecuación~\eqref{eq:endcorr} es
\emph{aproximada}; más adelante se mide $e$ de forma empírica con el método de regresión.

\subsection{Velocidad del sonido y temperatura}
La velocidad del sonido en aire seco depende de la temperatura absoluta. A partir de la
relación $\vsound=\sqrt{\gamma R T/M}$ se obtiene la forma práctica
\begin{equation}
  \vteo(T) \;=\; \SI{331.45}{\meter\per\second}\,\sqrt{1+\dfrac{T}{\SI{273.15}{\celsius}}},
  \label{eq:vteo}
\end{equation}
con $T$ en grados Celsius. Esta expresión es la que el programa usa como valor de referencia
(es exacta, a diferencia de la aproximación lineal $331.4+0.6\,T$).

\subsection{Relación fundamental}
Conocidos la frecuencia $f$ del diapasón y la longitud de onda $\lambda$ medida en la
resonancia, la velocidad experimental es
\begin{equation}
  \boxed{\;\vexp = f\,\lambda\;}
  \label{eq:fundamental}
\end{equation}

\section{Materiales}
\begin{itemize}[nosep]
  \item Tubo resonante vertical con depósito de agua de nivel ajustable.
  \item Cinta métrica adosada al tubo (resolución \SI{1}{\milli\meter}).
  \item Diapasón de frecuencia conocida (p.\,ej.\ \SI{440}{}, \SI{512}{} o \SI{1024}{\hertz}).
  \item Termómetro ambiente.
  \item Dispositivo con micrófono ejecutando el \emph{Asistente de Resonancia}.
\end{itemize}

\section{Procedimiento experimental paso a paso}

\subsection{Configuración inicial (panel ``Variables del Laboratorio'')}
\begin{enumerate}[nosep]
  \item Introducir la \textbf{frecuencia} $f$ grabada en el diapasón.
  \item Seleccionar el \textbf{armónico} a buscar (empezar por el Modo 1).
  \item Medir e introducir la \textbf{temperatura} $T$ del aire.
  \item Introducir el \textbf{diámetro} interno $D$ del tubo.
\end{enumerate}
El programa fija así la banda de análisis $f\pm\SI{15}{\hertz}$ y calcula $\vteo(T)$ y $e$.

\subsection{Predicción (panel ``Predictor de Resonancia'')}
Introducir el \textbf{largo del tubo}. El programa evalúa la
ecuación~\eqref{eq:resonancia} para todos los modos impares que caben y marca las longitudes
de columna de aire donde \emph{debe} oírse el máximo. Esto orienta la búsqueda antes de medir.

\subsection{Calibración del ruido}
Pulsar \emph{``Calibrar y Buscar Resonancia''} y guardar silencio \SI{3}{\second}. El sistema
promedia el ruido de fondo y fija un umbral de detección $u=\max(20,\ \bar{r}+15)$ en la escala
de amplitud (0--255), evitando falsos positivos por ruido ambiente.

\subsection{Captura de una resonancia}
\begin{enumerate}[nosep]
  \item Golpear el diapasón y acercarlo a la boca del tubo. El estado evoluciona
        \emph{Esperando $\to$ Estabilizando $\to$ Midiendo}; el filtro exige un tono
        sostenido durante varios fotogramas.
  \item Verificar que la \textbf{frecuencia detectada por la FFT} coincide con $f$
        (indicador en verde), lo que confirma que se mide el diapasón y no un ruido.
  \item Variar lentamente el nivel del agua. La \textbf{amplitud} y el \textbf{Nivel de
        Resonancia} crecen al acercarse al punto resonante.
  \item Sobrepasar ligeramente el máximo: cuando la amplitud cae de forma sostenida, el
        programa anuncia el pico.
  \item Leer en la cinta la \textbf{longitud de aire $L$} en el máximo e introducirla.
\end{enumerate}

\subsection{Repetición por modos}
Repetir la captura cambiando el armónico a Modo 3 y Modo 5, anotando $L_1$, $L_3$ y $L_5$
para el \emph{mismo} diapasón. Con dos modos basta para el método de regresión.

\section{Métodos de cálculo}

\subsection{Método directo (un punto)}
A partir de una sola resonancia de modo $m$ y longitud medida $L$:
\begin{align}
  L_{\text{ef}} &= L + e, & e &= 0.3\,D, \label{eq:Lef}\\
  \lambda &= \frac{4\,L_{\text{ef}}}{m}, &
  \vexp &= f\,\lambda. \label{eq:vexp_directo}
\end{align}
Su limitación es que depende de la corrección de extremo \emph{supuesta}
(ecuación~\eqref{eq:endcorr}).

\subsection{Método de diferencia de modos (regresión lineal)}\label{sec:regresion}
Es el método de mayor exactitud. Reescribiendo la ecuación~\eqref{eq:resonancia} como una recta
en la variable $m$:
\begin{equation}
  L_m = \underbrace{\frac{\lambda}{4}}_{\text{pendiente } a}\,m \;\underbrace{-\,e}_{\text{ordenada } b},
\end{equation}
un ajuste por mínimos cuadrados de los pares $(m,\,L_m)$ da
\begin{equation}
  a = \frac{\sum (m_i-\bar m)(L_i-\bar L)}{\sum (m_i-\bar m)^2},
  \qquad b = \bar L - a\,\bar m,
\end{equation}
y de ahí
\begin{equation}
  \boxed{\;\lambda = 4a,\qquad \vexp = 4 a f,\qquad e_{\text{medida}} = -\,b.\;}
  \label{eq:reg}
\end{equation}
\textbf{Ventaja clave:} la corrección de extremo se \emph{cancela} en la pendiente, por lo que
$\vexp$ no depende de ningún valor supuesto; además $e$ se \emph{mide} como
$-b$. El caso de dos modos es la conocida relación
\begin{equation}
  L_3 - L_1 = \frac{\lambda}{2}\;\Rightarrow\; \vexp = 2 f (L_3-L_1).
\end{equation}
La calidad del ajuste se cuantifica con el coeficiente de determinación
$R^2 = a^2\,S_{mm}/S_{LL}$, que debe ser muy próximo a 1.

\section{Análisis de incertidumbre}

\subsection{Método directo}
Aplicando propagación de errores a la ecuación~\eqref{eq:vexp_directo}:
\begin{equation}
  \frac{\Delta\vexp}{\vexp} = \sqrt{\left(\frac{\Delta f}{f}\right)^2
  + \left(\frac{\Delta L_{\text{ef}}}{L_{\text{ef}}}\right)^2},
  \label{eq:prop}
\end{equation}
con $\Delta L_{\text{ef}}=\Delta L$ la incertidumbre de lectura (por defecto
$\Delta L=\SI{0.2}{\centi\meter}$) y $\Delta f$ la tolerancia del diapasón
(por defecto $\Delta f=\SI{1}{\hertz}$).

\subsection{Método de regresión}
La incertidumbre de la velocidad proviene del error estándar de la pendiente:
\begin{equation}
  s_a = \sqrt{\frac{1}{n-2}\,\frac{\sum (L_i - a m_i - b)^2}{\sum (m_i-\bar m)^2}},
  \qquad \Delta\vexp = 4 f\, s_a .
  \label{eq:prop_reg}
\end{equation}
El resultado se reporta como $\vexp \pm \Delta\vexp$ y se considera \emph{compatible} con la
teoría si $|\vexp-\vteo|\le \Delta\vexp$.

\section{Ejemplo numérico resuelto ($f=\SI{1024}{\hertz}$)}

\paragraph{Datos.} $f=\SI{1024}{\hertz}$, $T=\SI{20}{\celsius}$, $D=\SI{3}{\centi\meter}$.
\begin{align*}
  \vteo &= 331.45\sqrt{1+20/273.15} = \SI{343.4}{\meter\per\second}, &
  \lambda &= \vteo/f = \SI{33.53}{\centi\meter},\\
  \lambda/4 &= \SI{8.38}{\centi\meter}, & e &= 0.3\cdot 3 = \SI{0.90}{\centi\meter}.
\end{align*}

\paragraph{Resonancias predichas (ec.~\eqref{eq:resonancia}).}
\begin{center}
\begin{tabular}{cccc}
\toprule
Modo $m$ & $m\lambda/4$ (cm) & $L_m = m\lambda/4 - e$ (cm) \\
\midrule
1 & 8.38  & 7.48 \\
3 & 25.15 & 24.25 \\
5 & 41.92 & 41.02 \\
\bottomrule
\end{tabular}
\end{center}

\paragraph{Medidas de laboratorio (ejemplo).}
$L_1=\SI{7.6}{\centi\meter}$, $L_3=\SI{24.3}{\centi\meter}$, $L_5=\SI{41.1}{\centi\meter}$.

\subparagraph{(a) Método directo, modo 1.}
\begin{align*}
  L_{\text{ef}} &= 7.6+0.9 = \SI{8.5}{\centi\meter}, &
  \lambda &= 4(0.085) = \SI{0.340}{\meter},\\
  \vexp &= 1024\cdot 0.340 = \SI{348.2}{\meter\per\second}, &
  \text{error} &= 1.4\,\%.
\end{align*}
Incertidumbre (ec.~\eqref{eq:prop}): $\Delta\vexp/\vexp=\sqrt{(1/1024)^2+(0.002/0.085)^2}=0.0236$,
luego $\vexp = \SI{348(8)}{\meter\per\second}$ (contiene a $\vteo$).

\subparagraph{(b) Diferencia de modos (1 y 3).}
\[
  \lambda = 2(L_3-L_1) = 2(0.243-0.076) = \SI{0.334}{\meter},\quad
  \vexp = 1024\cdot 0.334 = \SI{342.0}{\meter\per\second}\ (0.4\,\%).
\]

\subparagraph{(c) Regresión con tres modos.}
Con $m=\{1,3,5\}$ y $L=\{0.076,\,0.243,\,0.411\}\,\si{\meter}$:
\[
  a = \frac{0.670}{8} = \SI{0.08375}{\meter},\quad
  \lambda = 4a = \SI{0.335}{\meter},\quad
  \vexp = 4af = \SI{343.0}{\meter\per\second}.
\]
\[
  e_{\text{medida}} = -b = \SI{0.79}{\centi\meter},\qquad R^2 = 0.9999,\qquad
  \Delta\vexp = 4f s_a = \SI{0.3}{\meter\per\second}.
\]
\[
  \boxed{\;\vexp = \SI{343.0(3)}{\meter\per\second}\quad\text{frente a}\quad
  \vteo=\SI{343.4}{\meter\per\second}\;\;\Rightarrow\;\text{compatible}\ (0.1\,\%).\;}
\]
El método de regresión reduce el error de $1.4\,\%$ a $0.1\,\%$ y mide la corrección de
extremo ($\SI{0.79}{\centi\meter}$) sin suponerla.

\section{Fuentes de error y robustez del procedimiento}
Para que el resultado sea irreprochable, se acotan explícitamente todos los efectos:
\begin{itemize}[nosep]
  \item \textbf{Corrección de extremo:} eliminada en el método de regresión (no se supone, se mide).
  \item \textbf{Lectura de la longitud:} dominante; se propaga como $\Delta L=\SI{0.2}{\centi\meter}$.
        Su efecto se minimiza usando varios modos (mayor brazo en $m$).
  \item \textbf{Calibración del diapasón:} el diapasón es un \emph{patrón} calibrado; se usa su
        frecuencia nominal y la FFT la \emph{verifica} en vivo (banda $f\pm\SI{15}{\hertz}$).
  \item \textbf{Temperatura:} $\vteo$ varía $\approx\SI{0.6}{\meter\per\second}$ por
        \si{\celsius}; con $\Delta T=\SI{0.5}{\celsius}$ aporta $\Delta\vteo\approx\SI{0.3}{\meter\per\second}$.
  \item \textbf{Humedad relativa:} efecto de segundo orden ($\lesssim 0.4\,\%$ entre aire seco y
        saturado a \SI{20}{\celsius}); despreciable frente a la incertidumbre de lectura. Se
        asume aire seco (ec.~\eqref{eq:vteo}).
  \item \textbf{Pérdidas viscosas/térmicas en la pared:} afectan al factor de calidad (anchura
        del pico), no a la \emph{posición} de la resonancia, por lo que no sesgan $\vexp$.
  \item \textbf{Ruido ambiente y transitorios:} suprimidos por la calibración, el filtro de
        estabilidad y el suavizado exponencial de la señal.
\end{itemize}

\section{Conclusión}
El procedimiento determina $\vsound$ con dos métodos independientes y consistentes. El método de
diferencia de modos por regresión lineal proporciona el valor de referencia, pues elimina la
corrección de extremo, mide $e$ empíricamente y reporta una incertidumbre estadística rigurosa.
En el ejemplo, $\vexp=\SI{343.0(3)}{\meter\per\second}$ coincide con el valor teórico
$\vteo=\SI{343.4}{\meter\per\second}$ dentro de la incertidumbre, validando el experimento.

\appendix
\section{Correspondencia programa $\leftrightarrow$ ecuaciones}
\begin{center}
\begin{tabular}{ll}
\toprule
Elemento del programa & Ecuación / función \\
\midrule
``v teórica'' & \eqref{eq:vteo} \quad\texttt{theoreticalSpeedOfSound} \\
``Pico Encontrado'' $\to$ $\vexp\pm\Delta v$ & \eqref{eq:vexp_directo}, \eqref{eq:prop} \quad\texttt{singlePointResult} \\
``Predictor de Resonancia'' & \eqref{eq:resonancia} \quad\texttt{predictedResonances} \\
``Velocidad por Regresión'' & \eqref{eq:reg}, \eqref{eq:prop_reg} \quad\texttt{regressionResult} \\
``Frecuencia detectada (FFT)'' & interpolación parabólica del pico \\
\bottomrule
\end{tabular}
\end{center}

\end{document}
